% GENERATED FILE. Edit canonical lessons, then run npm run generate:books.
% canonical-source-hash: fnv1a64:7606f54a2b9b85fb
% canonical-lessons: GU-C436-lu, GU-C436-bapharo, GU-C436-paro, GU-C436-dajhadvu, GU-C436-sekavu

\chapter{Hot Winds and Mercury}
\label{ch:gu-hot-winds-and-mercury}
\hlchaptermodality{436}

\begin{chapteropening}
\textbf{By the end of this chapter:} \emph{I can talk about a very hot day.}

Complete the last lesson of chapter 436: I can talk about a very hot day.
\end{chapteropening}

\section[\texorpdfstring{\gu{લૂ} (lū)}{લૂ (lū)}]{\gu{લૂ} (lū) — a hot summer wind}

\label{lesson:GU-C436-lu}



\begin{quote}
Before the new one: say the Gujarati for a forecast, then the Gujarati for to flash.
\end{quote}

\subsection*{You'll want to know: \gu{લૂ}}
\textbf{\gu{લૂ}} — \emph{lū} — \textquotedblleft{}a hot summer wind\textquotedblright{}.

\begin{quote}
\gu{મે} \gu{મહિનામાં} \gu{બપોરે} \gu{ગરમ} \gu{લૂ} \gu{ફૂંકાય} \gu{છે}.
\end{quote}

\begin{quote}
\emph{me mahināmā̃ bapore garam lū phū̃kāy che.}
\end{quote}

\begin{quote}
In May a hot wind blows in the afternoon.
\end{quote}

A feminine noun.

\begin{cousinweb}
\textbf{\gu{લૂ}} (\emph{lū}) CDIAL lists Hindi \emph{lūhar} and Sindhi \emph{lūha}, \textquotedblleft{}the hot summer wind\textquotedblright{}, under \emph{*lūṣā}, \textquotedblleft{}burning\textquotedblright{}, a form it reconstructs (CDIAL 11099); it gives no Gujarati form, so the link is hedged.
\end{cousinweb}

\begin{grammarlens}[title={What you've built}]
One more word for getting around.
\end{grammarlens}

\subsection*{Your turn}
\begin{itemize}
\raggedright
  \item \emph{Say it:} \emph{lū}
  \item \emph{Say it:} \emph{lū}, once more
  \item \emph{Say it:} \emph{me mahināmā̃ bapore garam lū phū̃kāy che.} — In May a hot wind blows in the afternoon.
  \item \emph{Recall:} say the Gujarati for a forecast, then the Gujarati for to flash, then say \emph{lū} again
\end{itemize}

\begin{tcolorbox}[breakable,colback=teal!4,colframe=teal!35!black,title={Before you move on}]
What does \gu{લૂ} mean? (\textquotedblleft{}A hot summer wind\textquotedblright{}.) Say it once more.
\end{tcolorbox}
\section[\texorpdfstring{\gu{બફારો} (baphāro)}{બફારો (baphāro)}]{\gu{બફારો} (baphāro) — steamy heat (airless, before rain)}

\label{lesson:GU-C436-bapharo}



\begin{quote}
Before the new one: say the Gujarati for to flash, then the Gujarati for a hot summer wind.
\end{quote}

\subsection*{You'll want to know: \gu{બફારો}}
\textbf{\gu{બફારો}} — \emph{baphāro} — \textquotedblleft{}steamy heat (airless, before rain)\textquotedblright{}.

\begin{quote}
\gu{વરસાદ} \gu{પહેલાં} \gu{ઘણો} \gu{બફારો} \gu{હતો}.
\end{quote}

\begin{quote}
\emph{varsād pahelā̃ ghaṇo baphāro hato.}
\end{quote}

\begin{quote}
Before the rain it was very steamy.
\end{quote}

A masculine noun.

\begin{cousinweb}
\textbf{\gu{બફારો}} (\emph{baphāro}) goes back to \emph{*bāṣpākāra}, \textquotedblleft{}having the form of vapour\textquotedblright{}, a form CDIAL reconstructs; CDIAL lists the Gujarati word, \textquotedblleft{}vapour, sultriness\textquotedblright{} (CDIAL 9224).
\end{cousinweb}

\begin{grammarlens}[title={What you've built}]
One more word for getting around.
\end{grammarlens}

\subsection*{Your turn}
\begin{itemize}
\raggedright
  \item \emph{Say it:} \emph{baphāro}
  \item \emph{Say it:} \emph{baphāro}, once more
  \item \emph{Say it:} \emph{varsād pahelā̃ ghaṇo baphāro hato.} — Before the rain it was very steamy.
  \item \emph{Recall:} say the Gujarati for to flash, then the Gujarati for a hot summer wind, then say \emph{baphāro} again
\end{itemize}

\begin{tcolorbox}[breakable,colback=teal!4,colframe=teal!35!black,title={Before you move on}]
What does \gu{બફારો} mean? (\textquotedblleft{}Steamy heat (airless, before rain)\textquotedblright{}.) Say it once more.
\end{tcolorbox}
\section[\texorpdfstring{\gu{પારો} (pāro)}{પારો (pāro)}]{\gu{પારો} (pāro) — mercury; the temperature (on a thermometer)}

\label{lesson:GU-C436-paro}



\begin{quote}
Before the new one: say the Gujarati for a hot summer wind, then the Gujarati for steamy heat.
\end{quote}

\subsection*{You'll want to know: \gu{પારો}}
\textbf{\gu{પારો}} — \emph{pāro} — \textquotedblleft{}mercury; the temperature (on a thermometer)\textquotedblright{}.

\begin{quote}
\gu{આજે} \gu{પારો} \gu{ચાળીસ} \gu{ડિગ્રી} \gu{પર} \gu{પહોંચ્યો}.
\end{quote}

\begin{quote}
\emph{āje pāro cāḷīs ḍigrī par pahõcyo.}
\end{quote}

\begin{quote}
Today the temperature reached forty degrees.
\end{quote}

A masculine noun.

\begin{cousinweb}
\textbf{\gu{પારો}} (\emph{pāro}) is a tadbhava of Sanskrit \emph{pārada}, \textquotedblleft{}quicksilver\textquotedblright{}; CDIAL lists the Gujarati word (CDIAL 8104). The mercury in a thermometer gave the news its \gu{પારો}.
\end{cousinweb}

\begin{grammarlens}[title={What you've built}]
One more word for getting around.
\end{grammarlens}

\subsection*{Your turn}
\begin{itemize}
\raggedright
  \item \emph{Say it:} \emph{pāro}
  \item \emph{Say it:} \emph{pāro}, once more
  \item \emph{Say it:} \emph{āje pāro cāḷīs ḍigrī par pahõcyo.} — Today the temperature reached forty degrees.
  \item \emph{Recall:} say the Gujarati for a hot summer wind, then the Gujarati for steamy heat, then say \emph{pāro} again
\end{itemize}

\begin{tcolorbox}[breakable,colback=teal!4,colframe=teal!35!black,title={Before you move on}]
What does \gu{પારો} mean? (\textquotedblleft{}Mercury; the temperature (on a thermometer)\textquotedblright{}.) Say it once more.
\end{tcolorbox}
\section[\texorpdfstring{\gu{દઝાડવું} (dajhāḍvũ)}{દઝાડવું (dajhāḍvũ)}]{\gu{દઝાડવું} (dajhāḍvũ) — to scorch; to burn (someone)}

\label{lesson:GU-C436-dajhadvu}



\begin{quote}
Before the new one: say the Gujarati for steamy heat, then the Gujarati for mercury.
\end{quote}

\subsection*{You'll want to know: \gu{દઝાડવું}}
\textbf{\gu{દઝાડવું}} — \emph{dajhāḍvũ} — \textquotedblleft{}to scorch; to burn (someone)\textquotedblright{}.

\begin{quote}
\gu{બપોરનો} \gu{તડકો} \gu{પગને} \gu{દઝાડે} \gu{છે}.
\end{quote}

\begin{quote}
\emph{baporno taḍko pagne dajhāḍe che.}
\end{quote}

\begin{quote}
The afternoon sun scorches your feet.
\end{quote}

Past, with \gu{તેણે} (he or she …): \gu{તેણે} \gu{હાથ} \gu{દઝાડ્યો}, \emph{teṇe hāth dajhāḍyo}.

\begin{cousinweb}
\textbf{\gu{દઝાડવું}} (\emph{dajhāḍvũ}) is the causative partner of \gu{દાઝવું}, \textquotedblleft{}to be burnt\textquotedblright{}, which CDIAL lists under Sanskrit \emph{dahyatē}, \textquotedblleft{}is burnt\textquotedblright{} (CDIAL 6248): to make something burn.
\end{cousinweb}

\begin{grammarlens}[title={What you've built}]
One more word for getting around.
\end{grammarlens}

\subsection*{Your turn}
\begin{itemize}
\raggedright
  \item \emph{Say it:} \emph{dajhāḍvũ}
  \item \emph{Say it:} \emph{dajhāḍvũ}, once more
  \item \emph{Say it:} \emph{baporno taḍko pagne dajhāḍe che.} — The afternoon sun scorches your feet.
  \item \emph{Recall:} say the Gujarati for steamy heat, then the Gujarati for mercury, then say \emph{dajhāḍvũ} again
\end{itemize}

\begin{tcolorbox}[breakable,colback=teal!4,colframe=teal!35!black,title={Before you move on}]
What does \gu{દઝાડવું} mean? (\textquotedblleft{}To scorch; to burn (someone)\textquotedblright{}.) Say it once more.
\end{tcolorbox}
\section[\texorpdfstring{\gu{શેકાવું} (śekāvũ)}{શેકાવું (śekāvũ)}]{\gu{શેકાવું} (śekāvũ) — to bake (in the heat); to be roasted}

\label{lesson:GU-C436-sekavu}



\begin{quote}
Before the new one: say the Gujarati for mercury, then the Gujarati for to scorch.
\end{quote}

\subsection*{You'll want to know: \gu{શેકાવું}}
\textbf{\gu{શેકાવું}} — \emph{śekāvũ} — \textquotedblleft{}to bake (in the heat); to be roasted\textquotedblright{}.

\begin{quote}
\gu{આવા} \gu{તડકામાં} \gu{અમે} \gu{શેકાઈ} \gu{ગયા}.
\end{quote}

\begin{quote}
\emph{āvā taḍkāmā̃ ame śekāī gayā.}
\end{quote}

\begin{quote}
We got roasted in sunshine like this.
\end{quote}

Past (he …): \gu{તે} \gu{તડકામાં} \gu{શેકાયો}, \emph{te taḍkāmā̃ śekāyo}.

\begin{cousinweb}
\textbf{\gu{શેકાવું}} (\emph{śekāvũ}) is the intransitive partner of \gu{શેકવું}, \textquotedblleft{}to roast\textquotedblright{}, which goes back to \emph{*sēkk}, \textquotedblleft{}to heat, to roast\textquotedblright{}, a form CDIAL reconstructs (CDIAL 13581).
\end{cousinweb}

\begin{grammarlens}[title={What you've built}]
One more word for getting around.
\end{grammarlens}

\subsection*{Your turn}
\begin{itemize}
\raggedright
  \item \emph{Say it:} \emph{śekāvũ}
  \item \emph{Say it:} \emph{śekāvũ}, once more
  \item \emph{Say it:} \emph{āvā taḍkāmā̃ ame śekāī gayā.} — We got roasted in sunshine like this.
  \item \emph{Recall:} say the Gujarati for mercury, then the Gujarati for to scorch, then say \emph{śekāvũ} again
\end{itemize}

\begin{tcolorbox}[breakable,colback=teal!4,colframe=teal!35!black,title={Before you move on}]
What does \gu{શેકાવું} mean? (\textquotedblleft{}To bake (in the heat); to be roasted\textquotedblright{}.) Say it once more.
\end{tcolorbox}
